% XML Nodes auswählen
\startxmlsetups xml:jatssetups
	% Alles löschen und nur holen, was explizit ausgewählt wird
	\xmlsetsetup{#1}{*}{-}
	% Hauptgruppen
	\xmlsetsetup{#1}{article|front|body|back}{xml:*}
	% Appendix:
	\xmlsetsetup{#1}{app|app-group}{xml:*}
	\xmlsetsetup{#1}{app[@content-type='figures']}{xml:app-figures}
	% Metadaten
	\xmlsetsetup{#1}{article-meta}{xml:*}
	% Titelei
	\xmlsetsetup{#1}{contrib-group|contrib|name|string-name|surname|given-names|aff|email}{xml:*}
	\xmlsetsetup{#1}{front/notes}{xml:front:notes}
	% Sec
	\xmlsetsetup{#1}{sec[not @sec-type]}{xml:*}
	\xmlsetsetup{#1}{sec[@sec-type="nonum"]}{xml:sec-nonum}
	\xmlsetsetup{#1}{sec[@sec-type="thought-break"]}{xml:thought-break}
	\xmlsetsetup{#1}{sec[@sec-type="thought-break-asterism"]}{xml:thought-break-asterism}
	% Paragraphen
	\xmlsetsetup{#1}{p}{xml:*}
	\xmlsetsetup{#1}{preformat}{xml:preformat}
	% Blockquotes
	\xmlsetsetup{#1}{disp-quote}{xml:*}
	\xmlsetsetup{#1}{disp-quote[@content-type = 'blockquotenoindenting']}{xml:blockquotenoindenting}
	% Divs
	\xmlsetsetup{#1}{boxed-text[@content-type = 'dedication']}{xml:dedication}
	\xmlsetsetup{#1}{boxed-text[@content-type = 'epigraph']}{xml:epigraph}
	\xmlsetsetup{#1}{disp-quote[@content-type = 'epigraph']}{xml:epigraph}
	\xmlsetsetup{#1}{attrib}{xml:attrib}
	\xmlsetsetup{#1}{boxed-text[@content-type = 'signature']}{xml:signature}
	\xmlsetsetup{#1}{boxed-text[@content-type = 'lettersignature']}{xml:lettersignature}
	\xmlsetsetup{#1}{boxed-text[@content-type = 'flushleft']}{xml:flushleft}
	\xmlsetsetup{#1}{boxed-text[@content-type = 'flushright']}{xml:flushright}
	\xmlsetsetup{#1}{boxed-text[@content-type = 'centered']}{xml:centered}
	\xmlsetsetup{#1}{boxed-text[@content-type = 'righttoleft']}{xml:righttoleft}
	\xmlsetsetup{#1}{boxed-text[@content-type = 'blockquotenoindenting']}{xml:blockquotenoindenting}
	% Fussnoten
	\xmlsetsetup{#1}{xref[@ref-type='fn']}{xml:footnote}
    \xmlfilter{#1}{fn-group/fn/command(xml:collectnotes)}
	% Listen
	\xmlsetsetup{#1}{list[@list-type='bullet']}{xml:itemize}
	\xmlsetsetup{#1}{list[@list-type='order']}{xml:enumerate}
	\xmlsetsetup{#1}{list[@list-type='alpha-lower']}{xml:alphalist}
	\xmlsetsetup{#1}{list[@list-type='alpha-upper']}{xml:alphalist-uppercase}
	% Definition List
	\xmlsetsetup{#1}{def-list}{xml:*}
	\xmlsetsetup{#1}{def-list/title}{xml:def-list:title}
	\xmlsetsetup{#1}{def-list/term-head}{xml:def-list:term-head}
	\xmlsetsetup{#1}{def-list/def-head}{xml:def-list:def-head}
	\xmlsetsetup{#1}{def-list/def-item}{xml:def-list:def-item}
	\xmlsetsetup{#1}{def-list/def-item/term}{xml:def-list:def-item:term}
	\xmlsetsetup{#1}{def-list/def-item/def}{xml:def-list:def-item:def}
	% Formatierungen
	\xmlsetsetup{#1}{bold|italic|strike|underline|sc|styled-content|sup|sub}{xml:*}
	% Links
	\xmlsetsetup{#1}{ext-link}{xml:*}
	% Grafiken
	\xmlsetsetup{#1}{fig}{xml:*}
	\xmlsetsetup{#1}{fig/graphic}{xml:fig:*}
	\xmlsetsetup{#1}{graphic}{xml:*}
	% No graphic figures
	\xmlsetsetup{#1}{fig[@content-type='no-graphic']}{xml:no-graphic-figure}
	% Figure Group
	\xmlsetsetup{\xmldocument}{fig-group}{xml:*}
	\xmlsetsetup{\xmldocument}{fig-group/caption}{xml:figure-group:caption}
	\xmlsetsetup{\xmldocument}{fig-group/fig}{xml:figure-group:fig}
	\xmlsetsetup{\xmldocument}{fig-group/fig/caption}{xml:figure-group:fig:caption}
	\xmlsetsetup{\xmldocument}{fig-group/fig/graphic}{xml:figure-group:fig:graphic}
	% Etc
	\xmlsetsetup{#1}{xref[@ref-type='bibr']}{xml:citeref}
	% Bibliografie
	\xmlsetsetup{#1}{ref-list[@content-type="wrapper"]}{xml:ref-list-wrapper}
	\xmlsetsetup{#1}{ref-list[not @content-type]}{xml:ref-list}
	\xmlsetsetup{#1}{ref|mixed-citation}{xml:*}
	% Siglen
	\xmlsetsetup{#1}{ref-list[@content-type="shorthands"]}{xml:list-of-shorthands}
	\xmlsetsetup{#1}{ref-list[@content-type="shorthands"]/ref}{xml:shorthand}
	% Milestones
	\xmlsetsetup{#1}{milestone-start|milestone-end}{xml:*}
	% Named Content
	\xmlsetsetup{#1}{named-content}{xml:*}
	% Tabellen
	\xmlsetsetup{#1}{table-wrap[@content-type!='parallel' and @content-type!='parallel-4']}{xml:table-wrap}
	\xmlsetsetup{#1}{table[@content-type!='parallel' and @content-type!='parallel-r2l']}{xml:table}
	%\xmlsetsetup{#1}{thead|tbody|tr|th|td}{xml:*}
	% Paralleltexte
	\xmlsetsetup{#1}{boxed-text[@content-type='parallel-verse']}{xml:parallel-verse}
	\xmlsetsetup{#1}{table-wrap[@content-type='parallel']}{xml:table-wrap-parallel}
	\xmlsetsetup{#1}{table-wrap[@content-type='parallel-4']}{xml:table-wrap-parallel-4}
	\xmlsetsetup{#1}{table[@content-type='parallel']}{xml:parallel-text}
	\xmlsetsetup{#1}{table[@content-type='parallel-4']}{xml:parallel-text-4}
	\xmlsetsetup{#1}{table[@content-type='parallel-r2l']}{xml:parallel-text-r2l}
	% verse-groups
	\xmlsetsetup{#1}{verse-group[@content-type='verse-group-wrapper']}{xml:verse-group-wrapper}
	\xmlsetsetup{#1}{verse-group[@content-type='verse-group-wrapper-centered']}{xml:verse-group-wrapper-centered}
	\xmlsetsetup{#1}{verse-group[@content-type='stanza']}{xml:verse-group-stanza}
	\xmlsetsetup{#1}{verse-group[@content-type='stanza']/attrib}{xml:verse-group-stanza:attrib}
	\xmlsetsetup{#1}{verse-line}{xml:*}
	%\xmlstripanywhere{#1}{!pre}
	
	% Koltai addons
	\xmlsetsetup{#1}{fig[@content-type='koltai-figure']}{xml:koltai-figure}
	\xmlsetsetup{#1}{fig[@content-type='koltai-figure-table']}{xml:koltai-figure-table}
	\xmlsetsetup{#1}{table[@content-type='koltai-parallel']}{xml:koltai:parallel}
	\xmlsetsetup{#1}{table[@content-type='koltai-parallel-2']}{xml:koltai:parallel-text-2}
	\xmlsetsetup{#1}{fig[@content-type='koltai-figure']/boxed-text[@content-type='centered']}{xml:koltai-figure:boxed-text}
	\stopxmlsetups

\xmlregistersetup{xml:jatssetups}


%%%%%%%%%%%%%%%%%%%%%%%%%%%%%%%%%%%%%%%%%%%%%%%%%%%%%%%
%%%%%%%%%%%%%%%%%%%%%%%%%%%%%%%%%%%%%%%%%%%%%%%%%%%%%%%
% XML => ConTeXt
%%%%%%%%%%%%%%%%%%%%%%%%%%%%%%%%%%%%%%%%%%%%%%%%%%%%%%%
%%%%%%%%%%%%%%%%%%%%%%%%%%%%%%%%%%%%%%%%%%%%%%%%%%%%%%%

\def\xmltexdirective#1#2{\doif{#1}{command}{#2}}
\xmlinstalldirective{tex}{xmltexdirective}

% XML Entities

\xmltexentity {nbsp} {~}

% Injectors
\startsetups xml:directive:injector:noindent
  \noindentation
\stopsetups

\startsetups xml:directive:injector:setRL
  \setupalign[righttoleft]
\stopsetups

\startsetups xml:directive:injector:setLR
  \setupalign[lefttoright]
\stopsetups

\startsetups xml:directive:injector:judeoarabic
  \switchtobodyfont[hebrew]
\stopsetups

\startsetups xml:directive:injector:tolerant
  \setupalign[verytolerant,stretch]
\stopsetups
\startsetups xml:directive:injector:flushright
  
\stopsetups

\startsetups xml:directive:injector:newline
  \crlf
\stopsetups

\startsetups xml:directive:injector:linebreak
  %\break
  \penalty-10000
\stopsetups

\startsetups xml:directive:injector:lesslooseness
  \looseness=-1
\stopsetups

\startsetups xml:directive:injector:morelooseness
  \looseness=1
\stopsetups

\startsetups xml:directive:injector:addlinetopage
  \adaptlayout[lines=+1]
\stopsetups

\startsetups xml:directive:injector:bTablehead
  \bTABLEhead
\stopsetups

\startsetups xml:directive:injector:eTablehead
  \eTABLEhead
\stopsetups


% Befehle
%  \xmlfilter{#1}{/title/command(xml:title)}
%  \xmlflush{#1}

%%%%%%%%%%%%%%%%%%%%%%%%%%%%%%%%%%%%%%%%%%%%

% Sprache:
% Hier können wir diesen Mechanismus einsetzen:
% \xmlmapvalue{emph}{bold} {\bf}
% \xmlmapvalue{emph}{italic}{\it}
% Aufrufen so:
% \xmlvalue{emph}{\xmlatt{#1}{type}}{}

% Wir brauchen mainlanguage am Anfang; bei den übrigen Elementen braucht es language
% Beim Artikel schalten wir auf \mainlanguage[xml:lang] 
% Welche Elemente brauchen language? 
% Aktuell mal nur für einfache <p> => denn diese brauchen wir ja immer/meist für Textinhalte;
% Ausserdem für <boxed-text>, wird aktuell für divs von verwendet.

\xmlmapvalue{language}{de}{\language[de]}
\xmlmapvalue{language}{en}{\language[en]}
\xmlmapvalue{language}{fr}{\language[fr]}
\xmlmapvalue{language}{he}{\language[he]\switchtobodyfont[hebrew]}
\xmlmapvalue{language}{agr}{\language[agr]}

\startxmlsetups xml:language
	\xmlvalue{language}{\xmlatt{#1}{xml:lang}}{}
\stopxmlsetups

\startxmlsetups xml:mainlanguage
	%ATTRIBUT DA!
	%\xmlatt{#1}{xml:lang}
	\mainlanguage[\xmlatt{#1}{xml:lang}]
\stopxmlsetups

% Typ des Beitrags holen und entsprechenden Mode aktivieren / funktioniert nicht!
\startxmlsetups xml:article-type
	\enablemode[\xmlatt{#1}{article-type}]
\stopxmlsetups

% Hauptelement + Main Nodes

\startxmlsetups xml:article
	\begingroup
		% Type des Beitrags holen / funktioniert nicht
		\xmlfilter{#1}{.[@article-type]/command(xml:article-type)}
		% Sprache auswählen
		\xmlfilter{#1}{.[@xml:lang]/command(xml:mainlanguage)}
		% Restlicher Inhalt
		\xmlflush{#1}
	\endgroup
\stopxmlsetups

% Front
\startxmlsetups xml:front
	\xmlflush{#1}
\stopxmlsetups

% Body
\startxmlsetups xml:body
   	\directsetup{document:start}
	  \xmlflush{#1}
   	\directsetup{document:stop}
\stopxmlsetups

% Back
\startxmlsetups xml:back
	\xmlflush{#1}
\stopxmlsetups

% Appendix
\startxmlsetups xml:app
\setuphead[section,subsection][number=no]
\setuphead[section][page=yes]

	\startstructurelevel[title=\xmlfilter{#1}{/title/command(xml:app-title)}]
	\xmlflush{#1}
	\stopstructurelevel
\stopxmlsetups

% Appendix Figures
\startxmlsetups xml:app-figures
\setuphead[section,subsection][number=no]
\setuphead[section][page=yes]
\setupfloat[figure][default={force,top,bottom,page}]

	\startstructurelevel[title=\xmlfilter{#1}{/title/command(xml:app-title)}]
	\xmlflush{#1}
	\stopstructurelevel
\stopxmlsetups

\startxmlsetups xml:app-title
	\xmlflush{#1}
\stopxmlsetups

% Appendix Group
\startxmlsetups xml:app-group
	\xmlflush{#1}
\stopxmlsetups

%%%%%%%%%%%%%%%%%%%%%%%%%%%%%%%%%%%%%%%

% Darstellungen
\startxmlsetups xml:no-graphic-figure
	%\startpostponingnotes
    \startplacefigure[title=\xmlfilter{#1}{/caption/command(xml:caption)}]
	  \xmlflush{#1}
    \stopplacefigure
	%\stoppostponingnotes
\stopxmlsetups

%%%%%%%%%%%%%%%%%%%%%%%%%%%%%%%%%%%%%%%%
% Grafiken

%<fig>
%  <caption>Katze</caption>
%  <graphic mimetype="image" mime-subtype="jpeg" xlink:href="katze.jpg" xlink:title="" />
%</fig>

% Grafik ohne Labels
\startxmlsetups xml:graphic
	%\startpostponingnotes
    \startplacegraphic[title=\xmlfilter{#1}{/caption/command(xml:caption)}]
	  \externalfigure[\xmlatt{#1}{xlink:href}]
    \stopplacegraphic
	%\stoppostponingnotes
\stopxmlsetups

% Grafik mit Label
\startxmlsetups xml:fig
	%\startpostponingnotes
    \startplacefigure[title=\xmlfilter{#1}{/caption/command(xml:caption)}]
	  \xmlfilter{#1}{/graphic/command(xml:fig:graphic)}
    \stopplacefigure
	%\stoppostponingnotes
\stopxmlsetups

\startxmlsetups xml:caption
  \xmlflush{#1}
\stopxmlsetups

\startxmlsetups xml:fig:graphic
  \externalfigure[\xmlatt{#1}{xlink:href}]
\stopxmlsetups

%\placefigure [place] {My Caption} {\externalfigure[myone]}


% Demo:
%<resource type='figure'>
%<caption>A picture of a cow.</caption>
%<content><external file="cow.pdf"/></content>
%</resource>

%\xmlsetsetup{demo}{resource[@type='figure']}{xml:demo:figure}
%\xmlsetsetup{demo}{external}{xml:demo:*}

%\startxmlsetups xml:demo:figure
%\placefigure
%{\xmlfirst{#1}{/caption}}
%{\xmlfirst{#1}{/content}}
%\stopxmlsetups

%\startxmlsetups xml:demo:external
%\externalfigure[\xmlatt{#1}{file}]
%\stopxmlsetups

% Figure Groups
\startxmlsetups xml:fig-group
  \startplacefigure[location=top,title=\xmlfilter{#1}{/caption/command(xml:fig-group:caption)}]
    \startcombination[distance=2mm,after=,nx=2,ny=1,location=middle]
      \xmlfilter{#1}{/fig/command(xml:fig-group:fig)}
    \stopcombination
  \stopplacefigure
\stopxmlsetups

\startxmlsetups xml:fig-group:caption
  \xmlflush{#1}
\stopxmlsetups

\startxmlsetups xml:fig-group:fig
  \startcontent
    \xmldoif {#1} {/graphic} {
      \xmlfilter{#1}{/graphic/command(xml:fig-group:fig:graphic)}  [width=.4\textwidth]
    }
  \stopcontent
  \startcaption
      \xmlfilter{#1}{/caption/command(xml:fig-group:caption)}
  \stopcaption
\stopxmlsetups

\startxmlsetups xml:fig-group:fig:caption
  \xmlflush{#1}
\stopxmlsetups

\startxmlsetups xml:fig-group:fig:graphic
  \externalfigure[\xmlatt{#1}{xlink:href}]
\stopxmlsetups

%Metadaten
\startxmlsetups xml:article-meta
	\setupdocument [
		pub-year=\xmlfilter{#1}{/pub-date/year/command(xml:article-meta:pubdate:year)},
		volume=\xmlfilter{#1}{/volume/command(xml:article-meta:volume)},
		issue=\xmlfilter{#1}{/issue/command(xml:article-meta:issue)},
		issue-title=\xmlfilter{#1}{/issue-title/command(xml:article-meta:issue-title)},
		doi=\xmlfilter{#1}{/article-id[@pub-id-type='doi']/command(xml:article-meta:doi)},
		elocation-id=\xmlfilter{#1}{/elocation-id/command(xml:article-meta:elocation-id)},
		title=\xmlfilter{#1}{/title-group/article-title/command(xml:article-meta:title-group:article-title)},
		subtitle=\xmlfilter{#1}{/title-group/subtitle/command(xml:article-meta:title-group:subtitle)},
		author={\AuthorList},
		abstract=\xmlfilter{#1}{/abstract/command(xml:article-meta:abstract)},
	]
	\setupinteraction [
		title=\xmlfilter{#1}{/title-group/article-title/command(xml:article-meta:title-group:article-title)},
		author={\xmlfilter{#1}{/contrib-group/contrib/name/given-names/command(xml:given-names)}\space\xmlfilter{#1}{/contrib-group/contrib/name/surname/command(xml:surname)}},
	]
   \xmlflush{#1}
\stopxmlsetups

\startxmlsetups xml:article-meta:abstract
	\xmlflush{#1}
\stopxmlsetups

\startxmlsetups xml:article-meta:pubdate
	\xmlflush{#1}
\stopxmlsetups

\startxmlsetups xml:article-meta:pubdate:year
	\xmlflush{#1}
\stopxmlsetups

%Volume, Issue, Issue-Tite
\startxmlsetups xml:article-meta:volume
	\xmlflush{#1}
\stopxmlsetups

\startxmlsetups xml:article-meta:issue
	\xmlflush{#1}
\stopxmlsetups

\startxmlsetups xml:article-meta:issue-title
	\xmlflush{#1}
\stopxmlsetups

% elocation id
\startxmlsetups xml:article-meta:elocation-id
	\xmlflush{#1}
\stopxmlsetups

\startxmlsetups xml:article-meta:doi
	\xmlflush{#1}
\stopxmlsetups

% Titel
\startxmlsetups xml:article-meta:title-group:article-title
    \ignorespaces\xmlflush{#1}
\stopxmlsetups

\startxmlsetups xml:article-meta:title-group:subtitle
    \ignorespaces\xmlflush{#1}
\stopxmlsetups

% Autoren
\startxmlsetups xml:contrib-group
	\xmlflush{#1}
\stopxmlsetups

\startxmlsetups xml:contrib
  \defineauthor[\xmlfirst{#1}{/name/surname}\xmlfirst{#1}{/name/given-names}]
  [name={\xmlfirst{#1}{/name/string-name}},
  given-names={\xmlfirst{#1}{/name/given-names}},
  surname={\xmlfirst{#1}{/name/surname}},
  email={\xmlfirst{#1}{/email}},
  affiliation={\xmlfirst{#1}{/aff}},
  role={\xmlfilter{#1}{/role/command(xml:article-meta:contrib-group:contrib:role)}},
  ]
  \addtocommalist {\xmlfirst{#1}{/name/surname}\xmlfirst{#1}{/name/given-names}} \AuthorList
\stopxmlsetups

\startxmlsetups xml:name
	\xmlflush{#1}
\stopxmlsetups

\startxmlsetups xml:surname
	\xmlflush{#1}
\stopxmlsetups

\startxmlsetups xml:given-names
	\xmlflush{#1}
\stopxmlsetups

\startxmlsetups xml:string-name
	\xmlflush{#1}
\stopxmlsetups

\startxmlsetups xml:article-meta:contrib-group:contrib:role
	\xmlflush{#1}
\stopxmlsetups

\startxmlsetups xml:aff
	\xmlflush{#1}
\stopxmlsetups

\startxmlsetups xml:email
	\xmlflush{#1}
\stopxmlsetups

\startxmlsetups xml:front:notes
	\setupdocument [
		notes={\xmlflush{#1}},
	]
\stopxmlsetups
%%%%%%%%%%%%%%%%%%%%%%%%%%%%%%%%%%%%%%%%%%%5


% Textelemente

% blockstyles
\xmlmapvalue{content-type}{righttoleft}{\righttoleft}

% Reguläre Paragraphen
\startxmlsetups xml:p
	\begingroup
		% Sprache falls Attribut vorhanden
		\xmlfilter{#1}{.[@xml:lang]/command(xml:language)}
		% Style falls Attribut vorhanden
		\xmlvalue{content-type}{\xmlatt{#1}{content-type}}{}
		% Inhalt
		%\currentmainlanguage 
		%\currentlanguage
		\xmlflush{#1}
	\endgroup
	\par % Ist das schlau, oder soll man \par besser in der Gruppe platzieren?
\stopxmlsetups

% Preformat Paragraphen
\startxmlsetups xml:preformat
	\begingroup
		% Sprache falls Attribut vorhanden
		\xmlfilter{#1}{.[@xml:lang]/command(xml:language)}
		% Style falls Attribut vorhanden
		\xmlvalue{blockstyle}{\xmlatt{#1}{style}}{}
		% Inhalt
		\startmylines
			\xmlflushlinewise{#1}
		\stopmylines
	\endgroup
\stopxmlsetups

% Boxed Text (divs)

% dedication
\startxmlsetups xml:dedication
	\begingroup
		\setupalign[flushright]
		% Sprache falls Attribut vorhanden
		\xmlfilter{#1}{.[@xml:lang]/command(xml:language)} 
		% Inhalt
		\xmlflush{#1}
	\endgroup
\stopxmlsetups

% signature
\startxmlsetups xml:signature
	\blank
	\begingroup
		\setupalign[flushleft]
		\setupindenting[no]
		% Sprache falls Attribut vorhanden
		\xmlfilter{#1}{.[@xml:lang]/command(xml:language)} 
		% Inhalt
		\xmlflush{#1}
	\endgroup
\stopxmlsetups

% lettersignature
\startxmlsetups xml:lettersignature
	\blank[small]
	\begingroup
		\setupalign[flushleft]
		\setupindenting[no]
		% Sprache falls Attribut vorhanden
		\xmlfilter{#1}{.[@xml:lang]/command(xml:language)} 
		% Inhalt
		\xmlflush{#1}
	\endgroup
\stopxmlsetups

% Flushleft
\startxmlsetups xml:flushleft
	\begingroup
		\setupalign[flushleft]
		\setupindenting[no]
		% Sprache falls Attribut vorhanden
		\xmlfilter{#1}{.[@xml:lang]/command(xml:language)} 
		% Inhalt
		\xmlflush{#1}
	\endgroup
\stopxmlsetups

% Flushright
\startxmlsetups xml:flushright
	\begingroup
		\setupalign[flushright]
		\setupindenting[no]
		% Sprache falls Attribut vorhanden
		\xmlfilter{#1}{.[@xml:lang]/command(xml:language)} 
		% Inhalt
		\xmlflush{#1}
	\endgroup
\stopxmlsetups

% Centered
\startxmlsetups xml:centered
	\begingroup
		\setupalign[center]
		\setupindenting[no]
		% Sprache falls Attribut vorhanden
		\xmlfilter{#1}{.[@xml:lang]/command(xml:language)} 
		% Inhalt
		\xmlflush{#1}
	\endgroup
\stopxmlsetups



% Disp-Quote-Noindenting und Blockquotenoindenting
\startxmlsetups xml:blockquotenoindenting
	\startblockquote
	  \setupindenting[none]
	  \setupwhitespace[small]
      \xmlflush{#1}
	\stopblockquote
\stopxmlsetups

% Right to left

\startxmlsetups xml:righttoleft
	\begingroup
		\setupalign[r2l]
		%\setupindenting[no]
		% Sprache falls Attribut vorhanden
		\xmlfilter{#1}{.[@xml:lang]/command(xml:language)} 
		% Inhalt
		\xmlflush{#1}
	\endgroup
\stopxmlsetups

% epigraph
\startxmlsetups xml:epigraph
	\begingroup
		\setupnarrower[left=1.5\parindent]\startnarrower[left]\setupalign[flushright]
		% Sprache falls Attribut vorhanden
		\xmlfilter{#1}{.[@xml:lang]/command(xml:language)} 
		% Inhalt
		\xmlflush{#1}
		\stopnarrower\blank\noindentation
	\endgroup
\stopxmlsetups

\startxmlsetups xml:attrib
	\xmlflush{#1}
\stopxmlsetups

% Named Content 
% aktuell brauchen wir das nur für die Sprachwahl
\startxmlsetups xml:named-content
	\bgroup\xmlfilter{#1}{.[@xml:lang]/command(xml:language)}\xmlflush{#1}\egroup
\stopxmlsetups

% Milestones
% deprecated for language switching
% keeping for compatibility
\startxmlsetups xml:milestone-start
	\bgroup\xmlfilter{#1}{.[@xml:lang]/command(xml:language)}\xmlflush{#1}
\stopxmlsetups

\startxmlsetups xml:milestone-end
	\egroup
\stopxmlsetups

% Formatierungen

% Superscript
\startxmlsetups xml:sup
	\high{\xmlflush{#1}}
\stopxmlsetups

% Subscript
\startxmlsetups xml:sub
	\low{\xmlflush{#1}}
\stopxmlsetups

% Fett
\startxmlsetups xml:bold
	{\bf \xmlflush{#1}}
\stopxmlsetups

% Kursiv
\startxmlsetups xml:italic
	\emph{\xmlflush{#1}}
	%{\ignorespaces\em \xmlflush{#1}\removeunwantedspaces}
	% \xmlstrip{#1}{}
	%{\em \xmlflush{#1}}
\stopxmlsetups

\startxmlsetups xml:sc
	{\sc \xmlflush{#1}}
\stopxmlsetups

% Durchgestrichen
\startxmlsetups xml:strike
	\overstrike{\xmlflush{#1}}
\stopxmlsetups

% Unterstrichen
\startxmlsetups xml:underline
	\underbar{\xmlflush{#1}}
\stopxmlsetups

% Unterstrichen mit Dots
\startxmlsetups xml:underdot
	\underdot{\xmlflush{#1}}
\stopxmlsetups


% Unterstrichen mit Dashes
\startxmlsetups xml:underdash
	\underdash{\xmlflush{#1}}
\stopxmlsetups


% Named content

% general mappings
\xmlmapvalue{named-content}{specialemph}{\inframed[background=color,backgroundcolor=gray,frame=off]}
\xmlmapvalue{named-content}{underdot}{\underdot}

\xmlmapvalue{named-content}{italics-underlined}{\italicsunderlined}
\define[1]\italicsunderlined{\underbar{\emph{#1}}}

\xmlmapvalue{named-content}{emph-red}{\emphred}
\define[1]\emphred{\color[middlered]{#1}}


% cb:textvar mappings to tex macros
\xmlmapvalue{named-content}{cb-textvarA}{\cbtextvarA}
\xmlmapvalue{named-content}{cb-textvarB}{\cbtextvarB}
\xmlmapvalue{named-content}{cb-textvarC}{\cbtextvarC}
\xmlmapvalue{named-content}{cb-textvarAB}{\cbtextvarAB}
\xmlmapvalue{named-content}{cb-textvarAC}{\cbtextvarAC}
\xmlmapvalue{named-content}{cb-textvarBC}{\cbtextvarBC}
\xmlmapvalue{named-content}{cb-textvarABC}{\cbtextvarABC}
% cb textvariants
\define[1]\cbtextvarA{\underbar{#1}}
% \define[1]\cbtextvarB{\underdot{#1}}
\define[1]\cbtextvarB{\inframed[background=color,backgroundcolor=gray,frame=off]{#1}}
%\define[1]\cbtextvarC{\emph{#1}}
\define[1]\cbtextvarC{\color[middlered]{#1}}
% cb textvariant combinations
\define[1]\cbtextvarAB{\cbtextvarA{\cbtextvarB{#1}}}
\define[1]\cbtextvarAC{\cbtextvarA{\cbtextvarC{#1}}}
\define[1]\cbtextvarBC{\cbtextvarB{\cbtextvarC{#1}}}
\define[1]\cbtextvarABC{\cbtextvarA{\cbtextvarB{\cbtextvarC{#1}}}}



% Bidi Overrides
\xmlmapvalue{named-content}{bdoltr}{\lefttoright}
\xmlmapvalue{named-content}{bdortl}{\righttoleft}



\startxmlsetups xml:named-content
	\xmlvalue{named-content}{\xmlatt{#1}{content-type}}{not found}{\xmlflush{#1}}
\stopxmlsetups


% Styled content

\definestartstop[LTR] [before={\begingroup\lefttoright},after=\endgroup]
\definestartstop[RTL] [before={\begingroup\righttoleft},after=\endgroup]

\xmlmapvalue{style}{direction: ltr}{\lefttoright}
\xmlmapvalue{style}{direction: rtl}{\righttoleft}
\xmlmapvalue{style}{hebrew}{\switchtobodyfont[hebrew]}

\startxmlsetups xml:styled-content
	{
	\xmldoif {#1} {.[@style]}{\xmlvalue{style}{\xmlatt{#1}{style}}{not found}}
	\xmlfilter{#1}{.[@xml:lang]/command(xml:language)}
	\xmlflush{#1}
	}
\stopxmlsetups


% Links

% \startxmlsetups xml:ext-link
    % \begingroup
        % \expandUx
% \expanded{\goto{\hyphenatedurl{\xmlflush{#1}}}[url(\xmlflush{#1})]}
    % \endgroup
% \stopxmlsetups 

% \startxmlsetups xml:ext-link
% \goto{\hyphenatedurl{\xmlflush{#1}}}[url(\xmlpure{#1})]
% \stopxmlsetups 

\startxmlsetups xml:ext-link
\goto{\hyphenatedurl{\xmlflush{#1}}}[url(\xmlflushpure{#1})]
\stopxmlsetups

%Thought Break
\usemodule[fancybreak]
\setupfancybreak[
	symbol=star,
	spacebefore=halfline,
	spaceafter=halfline,
	indentnext=no]

\startxmlsetups xml:thought-break
  \blank\noindentation
  \xmlflush{#1}
\stopxmlsetups

\startxmlsetups xml:thought-break-asterism
  \fancybreak
  \xmlflush{#1}
\stopxmlsetups

% Gliederungselemente
\startxmlsetups xml:sec
	\startstructurelevel [title=\xmlfilter{#1}{/title/command(xml:sec:title)}]
	 \xmlflush{#1}
	\stopstructurelevel
\stopxmlsetups

\startxmlsetups xml:sec-nonum
	\startstructurelevel [title=\xmlfilter{#1}{/title/command(xml:sec:title)},number=no]
	 \xmlflush{#1}
	\stopstructurelevel
\stopxmlsetups

\startxmlsetups xml:sec:title
	\xmlflush{#1}
\stopxmlsetups

% Fussnoten
% Wir verwenden die Lösung aus dem Handbuch (XML Mkiv) S. 86f.

\startluacode
userdata.notes = {}
\stopluacode

\startxmlsetups xml:collectnotes
\ctxlua{userdata.notes['\xmlatt{#1}{id}'] = '#1'}
\stopxmlsetups

\startxmlsetups xml:footnote
\footnote
{\xmlflush
{\cldcontext{userdata.notes['\xmlatt{#1}{rid}']}}}
\stopxmlsetups


\xmlmapvalue{dispquotestartvar}{direction: rtl}{\startrtlblockquote}
\xmlmapvalue{dispquotestopvar}{direction: rtl}{\stoptrtlblockquote}
\xmlmapvalue{dispquotestopvar}{direction: ltr}{\startblockquote}
\xmlmapvalue{dispquotestopvar}{direction: ltr}{\stopblockquote}


% Eingerückte Zitate
% \startxmlsetups xml:disp-quote
	% \startblockquote
	% \xmlflush{#1}
	% \stopblockquote
% \stopxmlsetups

\startxmlsetups xml:disp-quote
	\xmldoifelse {#1} {.[@style="direction: rtl"]} {
		\startrtlblockquote
		\xmlflush{#1}
		\stoprtlblockquote
	}
	{
		\startblockquote
		\xmlflush{#1}
		\stopblockquote
	}	
\stopxmlsetups

% Itemize
\startxmlsetups xml:itemize
	\startitemize
		\xmlfilter{#1}{/list-item/command(xml:itemize:list-item)}
	\stopitemize
\stopxmlsetups

\startxmlsetups xml:enumerate
	\startitemize[n]
		\xmlfilter{#1}{/list-item/command(xml:itemize:list-item)}
	\stopitemize
\stopxmlsetups

\startxmlsetups xml:alphalist
	\startitemize[a][stopper={)}]
		\xmlfilter{#1}{/list-item/command(xml:itemize:list-item)}
	\stopitemize
\stopxmlsetups

\startxmlsetups xml:alphalist-uppercase
	\startitemize[A][stopper={.}]
		\xmlfilter{#1}{/list-item/command(xml:itemize:list-item)}
	\stopitemize
\stopxmlsetups

\startxmlsetups xml:itemize:list-item
	\startitem
		\xmlflush{#1}
	\stopitem
\stopxmlsetups

% description lists

\startxmlsetups xml:def-list
	\xmlflush{#1}
\stopxmlsetups

\startxmlsetups xml:def-list:title
	\xmlflush{#1}
\stopxmlsetups

\startxmlsetups xml:def-list:term-head
	\xmlflush{#1}
\stopxmlsetups

\startxmlsetups xml:def-list:def-head
	\xmlflush{#1}
\stopxmlsetups

\startxmlsetups xml:def-list:def-item
	\xmlflush{#1}
\stopxmlsetups

\startxmlsetups xml:def-list:def-item:term
	\startdeflist{\xmlflush{#1}}
\stopxmlsetups

\startxmlsetups xml:def-list:def-item:def
	\xmlflush{#1}
	\stopdeflist
\stopxmlsetups

\definedescription[deflistold][
  headstyle=normal, 
  style=normal, 
  alternative=hanging,
  %hang=margin,
  width=fit, 
  %right={--},
  margin=1cm,
  distance=0em,
  ]
  
\definedescription
  [deflist]
  [alternative=serried,
   %hang=margin,
   width=fit,
   headstyle=normal, 
   style=normal, 
   headcommand=\groupedcommand{}{:},
   distance=.5em,
   margin=1cm,
   indentnext=no,
   ]
   
\definedescription
  [deflisttop]
  [alternative=top,
   width=fit,
   distance=1em,
   margin=1.5em,
   headstyle=normal,
   inbetween={\blank[small]},
   headcommand=\groupedcommand{}{:},
   indentnext=no]


% Verse-group

\startxmlsetups xml:verse-group-wrapper
	\startverse
		\xmlflush{#1}
	\stopverse
\stopxmlsetups

\startxmlsetups xml:verse-group-wrapper-centered
	\startcenteredverse
		\xmlflush{#1}
	\stopcenteredverse
\stopxmlsetups

\startxmlsetups xml:verse-group-stanza
	\xmlflush{#1}\blank[preference,standard]
\stopxmlsetups

\startxmlsetups xml:verse-line
	\xmlflush{#1}\blank[none]
\stopxmlsetups

\startxmlsetups xml:verse-group-stanza:attrib
	\blank[force,quarterline]{\tx \xmlflush{#1}}
\stopxmlsetups

% Paralleltexte

\xmlmapvalue{colstyle}{direction: rtl}{,righttoleft}
\xmlmapvalue{colstyle}{direction: ltr}{,lefttoright}

\startxmlsetups xml:parallel-table:column-def
	\xmlvalue{colstyle}{\xmlatt{#1}{style}}{}
\stopxmlsetups

\startxmlsetups xml:table-wrap-parallel
\xmldoif{#1}{/caption}{
	\paratextheading{\xmlfilter{#1}{/caption/command(xml:caption)}}}
\xmlfilter{#1}{/table/command(xml:parallel-text)}
\noindentation
\stopxmlsetups

\startxmlsetups xml:parallel-text
    \startpostponingnotes
	\starttabulate
			[|
			p(.48\textwidth)A{verytolerant,extremestretch\xmlfilter {#1} {/colgroup/col[position()=1]/command(xml:parallel-table:column-def)}}
			|
			p(.48\textwidth)A{verytolerant,extremestretch\xmlfilter {#1} {/colgroup/col[position()=2]/command(xml:parallel-table:column-def)}}
			|]
		\xmlfilter{#1}{/thead/command(xml:parallel-text:thead)}
		\xmlfilter{#1}{/tbody/command(xml:parallel-text:tbody)}
	\stoptabulate
	\stoppostponingnotes
\stopxmlsetups

\startxmlsetups xml:table-wrap-parallel-4
\xmldoif{#1}{/caption}{
	\paratextheading{\xmlfilter{#1}{/caption/command(xml:caption)}}}
\xmlfilter{#1}{/table/command(xml:parallel-text-4)}
\noindentation
\stopxmlsetups

\startxmlsetups xml:parallel-text-4
    \startpostponingnotes
	\starttabulate
			[|
			p(.23\textwidth)A{flushleft,verytolerant,extremestretch\xmlfilter {#1} {/colgroup/col[position()=1]/command(xml:parallel-table:column-def)}}
			|
			p(.23\textwidth)A{flushleft,verytolerant,extremestretch\xmlfilter {#1} {/colgroup/col[position()=2]/command(xml:parallel-table:column-def)}}
			|
			p(.23\textwidth)A{flushleft,verytolerant,extremestretch\xmlfilter {#1} {/colgroup/col[position()=3]/command(xml:parallel-table:column-def)}}
			|
			p(.23\textwidth)A{flushleft,verytolerant,extremestretch\xmlfilter {#1} {/colgroup/col[position()=4]/command(xml:parallel-table:column-def)}}
			|]
		\xmlfilter{#1}{/thead/command(xml:parallel-text:thead)}
		\xmlfilter{#1}{/tbody/command(xml:parallel-text:tbody)}
	\stoptabulate
	\stoppostponingnotes
\stopxmlsetups

\startxmlsetups xml:koltai:parallel-text-2
	\startpostponingnotes
	\starttabulate
			[|
			p(.38\textwidth)A{verytolerant,extremestretch\xmlfilter {#1} {/colgroup/col[position()=1]/command(xml:parallel-table:column-def)}}
			|
			p(.25\textwidth)A{verytolerant,extremestretch\xmlfilter {#1} {/colgroup/col[position()=2]/command(xml:parallel-table:column-def)}}
			|]
		\xmlfilter{#1}{/thead/command(xml:parallel-text:thead)}
		\xmlfilter{#1}{/tbody/command(xml:parallel-text:tbody)}
	\stoptabulate
	\stoppostponingnotes
\stopxmlsetups

\startxmlsetups xml:parallel-text-r2l
    \startpostponingnotes
	\starttabulate[|pA{righttoleft,verytolerant,extremestretch}|pA{verytolerant,extremestretch}|]
		%\xmlfilter{#1}{/tr/command(xml:parallel-text:tr)}
		\xmlfilter{#1}{/thead/command(xml:parallel-text:thead)}
		\xmlfilter{#1}{/tbody/command(xml:parallel-text:tbody)}
	\stoptabulate
	\stoppostponingnotes
\stopxmlsetups

\startxmlsetups xml:parallel-text:tr
	\xmlfilter{#1}{/th/command(xml:parallel-text:tr:th)}
	\xmlfilter{#1}{/td/command(xml:parallel-text:tr:td)}
	\NC\NR
\stopxmlsetups

\startxmlsetups xml:parallel-text:tr:th
	\NC \xmlflush{#1}
\stopxmlsetups

\startxmlsetups xml:parallel-text:tr:td
	\NC \xmlflush{#1}
\stopxmlsetups

\startxmlsetups xml:parallel-text:thead:tr:th
	\NC \xmlflush{#1}
\stopxmlsetups


\startxmlsetups xml:parallel-text:thead
	\xmlfilter{#1}{/tr/command(xml:parallel-text:thead:tr)}
\stopxmlsetups

\startxmlsetups xml:parallel-text:thead:tr
	\xmlfilter{#1}{/th/command(xml:parallel-text:thead:tr:th)}\NC\NR
\stopxmlsetups

\startxmlsetups xml:parallel-text:thead:tr:th
	\NC \xmlflush{#1}
\stopxmlsetups

\startxmlsetups xml:parallel-text:tbody
	\xmlfilter{#1}{/tr/command(xml:parallel-text:tbody:tr)}
\stopxmlsetups

\startxmlsetups xml:parallel-text:tbody:tr
	\xmlfilter{#1}{/td/command(xml:parallel-text:tbody:tr:td)}\NC\NR
\stopxmlsetups

\startxmlsetups xml:parallel-text:tbody:tr:td
	\NC \xmlflush{#1}
\stopxmlsetups



% Parallel Verse
\defineparagraphs[paraverse]
				 [n=2,
				 %after={\noindentation},
				 distance=2em,
				 ]
				 
\setupparagraphs[paraverse][1][
   width=.45\textwidth,
   %distance=4em,
   align={verytolerant,stretch}]
   
\setupparagraphs[paraverse][2][
   %distance=4em,
   align={verytolerant,stretch}]
% width of 2nd column is calculated automatically

\startxmlsetups xml:parallel-verse
	\startparaverse
	%\xmlfirst{#1}{verse-group}
	\xmlfilter{#1}{/table-wrap/table/tbody/tr/td[position()=1]/command(xml:parallel-verse:col)}
	\paraverse
	%\nextparaverse
	%\xmllast{#1}{verse-group}
	 \xmlfilter{#1}{/table-wrap/table/tbody/tr/td[position()=2]/command(xml:parallel-verse:col)}
	\stopparaverse
\stopxmlsetups

\startxmlsetups xml:parallel-verse:col
	\xmlflush{#1}
\stopxmlsetups

% Tabellen

% \startxmlsetups xml:table-wrap
    % \xmldoif {#1} {.[@position="float"]}
	% {\startplacetable[location={split,none}]
	  % \begingroup
	% %\startalignment[middle]
	  % \blank[medium]\noindentation\placefloatcaption[table][title=\xmlfilter{#1}{/caption/command(xml:caption)}]
	% %\stopalignment
	% \endgroup
	% }
	  % \startpostponingnotes
	    % \xmlflush{#1}
	  % \stoppostponingnotes
    % \xmldoif {#1} {.[@position="float"]}
    % {\stopplacetable}
% \stopxmlsetups

\startxmlsetups xml:table-wrap
    \xmldoif {#1} {.[@position="float"]}
	{\startplacetable[location={split},title=\xmlfilter{#1}{/caption/command(xml:caption)}]}
	  \startpostponingnotes
	    \xmlflush{#1}
	  \stoppostponingnotes
    \xmldoif {#1} {.[@position="float"]}
    {\stopplacetable}
\stopxmlsetups

% \startxmlsetups xml:table-wrap
% \startlocalfootnotes
% \placetable
    % %The caption
    % {\xmlfilter{#1}{/caption/command(xml:caption)}}
    % %the table
    % \placelegend
        % {\xmlflush{#1}}
        % %the footnotes
        % {\placelocalfootnotes}
% \stoplocalfootnotes
% \stopxmlsetups


\xmlmapvalue{colgroupalign}{direction: rtl}{align=righttoleft}
\xmlmapvalue{colgroupalign}{direction: ltr}{align=lefttoright}

% \xmlmapvalue{colgroupalign}{direction: rtl}{righttoleft}
% \xmlmapvalue{colgroupalign}{direction: ltr}{lefttoright}

\startxmlsetups xml:table:colgroup
	\xmlfilter {#1} {/col/command(xml:table:colgroup:col)}
\stopxmlsetups

\startxmlsetups xml:table:colgroup:col
     \normalexpanded {
         \setupTABLE
         [c]
         [\xmlmatch {#1}]
         [\xmlvalue{colgroupalign}{\xmlatt{#1}{style}}{align={lefttoright,flushleft,verytolerant,extremestretch}},% align
		 ]
     }
\stopxmlsetups

\startxmlsetups xml:table
  \start
  \xmlfilter {#1} {/colgroup/command(xml:table:colgroup)}
  \bTABLE 
	[
	setups={table-teststyle},
	%offset=2mm,
	%spaceinbetween=5cm,
	%distance=1cm,
	]
	\xmlfilter{#1}{/thead/command(xml:table:thead)}
	\xmlfilter{#1}{/tbody/command(xml:table:tbody)}
  \eTABLE
  \stop
\stopxmlsetups

\startxmlsetups xml:table:thead
	\bTABLEhead
		\xmlfilter{#1}{/tr/command(xml:table:thead:tr)}
	\eTABLEhead
\stopxmlsetups

\startxmlsetups xml:table:thead:tr
	\bTR
		\xmlfilter{#1}{/th/command(xml:table:thead:tr:th)}
	\eTR
\stopxmlsetups

\startxmlsetups xml:table:thead:tr:th
	\bTH
		\xmlflush{#1}
	\eTH
\stopxmlsetups

\startxmlsetups xml:table:tbody 
\bTABLEbody\xmlfilter{#1}{/tr/command(xml:table:tbody:tr)}\eTABLEbody
\stopxmlsetups

\startxmlsetups xml:table:tbody:tr
\bTR\xmlfilter{#1}{/td/command(xml:table:tbody:tr:td)}\eTR
\stopxmlsetups


\startxmlsetups xml:table:tbody:tr:td
	\normalexpanded {
		\bTD [\xmlvalue{colgroupalign}{\xmlatt{#1}{style}}{}]
			\xmlflush{#1} 
		\eTD
	}
\stopxmlsetups

% Zitatnachweise
\startxmlsetups xml:citeref
	\xmlflush{#1}
\stopxmlsetups

% Bibliografie

\startxmlsetups xml:ref-list-wrapper
	\startstructurelevel [title={\xmlfilter{#1}{/title/command(xml:ref-list:title)}},number=no]
	 \xmlflush{#1}
 	\stopstructurelevel
\stopxmlsetups

\startxmlsetups xml:ref-list
	\startreferencelist
	\startstructurelevel [title={\xmlfilter{#1}{/title/command(xml:ref-list:title)}},number=no]
	 \xmlflush{#1}
 	\stopstructurelevel
	\stopreferencelist
\stopxmlsetups

\startxmlsetups xml:ref-list:title
	\xmlflush{#1}
\stopxmlsetups

\startxmlsetups xml:ref
	\xmlflush{#1}
\stopxmlsetups 

\startxmlsetups xml:mixed-citation
	\xmlflush{#1}\par
\stopxmlsetups 

% Siglen
\startxmlsetups xml:list-of-shorthands
	\startstructurelevel [title={\xmlfilter{#1}{/title/command(xml:ref-list:title)}},number=no]
		\starttabulate[|p(.15\textwidth)|pA{verytolerant,extremestretch}|] % TODO: LAYOUT prüfen!
			\xmlflush{#1}
		\stoptabulate
	\stopstructurelevel
\stoptabulate
\stopxmlsetups

\startxmlsetups xml:shorthand
  \xmlfilter{#1}{/label/command(xml:shorthand:label)} 
  \xmlfilter{#1}{/mixed-citation/command(xml:shorthand:content)} 
  \NC\NR
\stopxmlsetups

\startxmlsetups xml:shorthand:label
  \NC \xmlflush{#1}
\stopxmlsetups

\startxmlsetups xml:shorthand:content
  \NC \xmlflush{#1}
\stopxmlsetups

%\starttabulate[|rw(.25\textwidth)|lw(.75\textwidth)|] -> noch anpassen

% \starttabulate[|p(.1\textwidth)|p(.75\textwidth)|]
% \NC ND \NC Adorno, Negative Dialektik \NC\NR
% \NC label \NC mixed citation \NC\NR
% \stoptabulate

% siehe oben:
% \startxmlsetups xml:parallel-text
    % \startpostponingnotes
	% \starttabulate
			% [|
			% pA{verytolerant,extremestretch\xmlfilter {#1} {/colgroup/col[position()=1]/command(xml:parallel-table:column-def)}}
			% |
			% pA{verytolerant,extremestretch\xmlfilter {#1} {/colgroup/col[position()=2]/command(xml:parallel-table:column-def)}}
			% |]
		% \xmlfilter{#1}{/thead/command(xml:parallel-text:thead)}
		% \xmlfilter{#1}{/tbody/command(xml:parallel-text:tbody)}
	% \stoptabulate
	% \stoppostponingnotes
% \stopxmlsetups

%%%%%%%%%%%%%%%%%%%%%%%%%%%%%%%%%%%%%%%%%%%%%%%%%%%%%%%%%%%%%
% Koltai Hacks
% Darstellungen
\startxmlsetups xml:koltai-figure
	%\startpostponingnotes
    \startplacefigure[title=\xmlfilter{#1}{/caption/command(xml:caption)},location={page}]
	  \xmlflush{#1}
    \stopplacefigure
	%\stoppostponingnotes
\stopxmlsetups

\startxmlsetups xml:koltai-figure-table
	\start
	\startpostponingnotes
	\setupcaption
	  [table]
	  [
	   headstyle=smallcaps,
	   numberstopper={:},
	   suffix=,
	   prefix=yes,
	   number=yes,
	   distance=0em,
	   %way=bysection,
	   %prefixsegments=section,
	   %align=right,
	   location=bottom]
    \startplacefigure[title=\xmlfilter{#1}{/caption/command(xml:caption)},location={force,split}]
	  \xmlflush{#1}
    \stopplacefigure
	\stoppostponingnotes
	\stop

\stopxmlsetups

\startxmlsetups xml:koltai-figure:boxed-text
	\startalignment[center]
	  \xmlflush{#1}
	\stopalignment
\stopxmlsetups

% \startplacetable[number=no,location={force,split}]
% \startxtable
	% [
	% frame=off,
	% width=.3\textwidth,
	% ]
                  % \startxrow
                     % \startxcell In Israel\stopxcell
                     % \startxcell[loffset=10pt] when\stopxcell
                  % \stopxrow
% \stopxtable
% \stopplacetable

\definehspace[twoem][2 em]

\startxmlsetups xml:koltai:parallel
    \startembeddedxtable[frame=off,width=.3\textwidth,offset=0pt,split=yes]
	  \xmlfilter{#1}{/tbody/command(xml:koltai:parallel:tbody)}
	\stopembeddedxtable
\stopxmlsetups

% \startxmlsetups xml:koltai:parallel
    % \bTABLE[frame=off,width=.3\textwidth]
	  % \xmlfilter{#1}{/tbody/command(xml:table:tbody)}
	% \eTABLE
% \stopxmlsetups

\startxmlsetups xml:koltai:parallel:tbody
	\xmlfilter{#1}{/tr/command(xml:koltai:parallel:tbody:tr)}
\stopxmlsetups

\startxmlsetups xml:koltai:parallel:tbody:tr
	\startxrow 
	  \xmlfilter{#1}{/td/command(xml:koltai:parallel:tbody:tr:td)} 
	\stopxrow
\stopxmlsetups


\startxmlsetups xml:koltai:parallel:tbody:tr:td
\doifnumberelse
	{\xmlatt{#1}{colspan}}
	{\startxcell [nx=\xmlatt{#1}{colspan},align=middle] \xmlflush{#1} \stopxcell}
	{\startxcell \xmlflush{#1} \stopxcell}
\stopxmlsetups


